% Vietnamese translation for OI Public Library
% Source-File: IOI中国国家候选队论文/2005/何林/何林.doc
\documentclass[11pt]{article}
\usepackage{oipl-vn}

\title{Đơn giản hóa quan hệ dữ liệu}
\author{He Lin}
\date{2005}

\begin{document}
\maketitle

\origtitle{Simplifying data relations}
\sourcepath{IOI China candidate team papers / 2005 / He Lin / He Lin.doc}

\begin{abstract}
Quan hệ giữa dữ liệu quan trọng không kém bản thân dữ liệu. Các quan hệ thường
gặp nhất là quan hệ tuyến tính, quan hệ cây và quan hệ đồ thị. Bản chất của
thi tin học là khai thác đầy đủ dữ liệu, nhưng đôi khi khai thác quá đầy đủ
lại làm bài toán phức tạp hơn. Bài viết này bàn về cách đơn giản hóa quan hệ
dữ liệu một cách thích hợp để khai thác dữ liệu hợp lý.
\end{abstract}

\noindent\textbf{Từ khóa:} cây, đồ thị, dãy, quan hệ dữ liệu.

\section{Dẫn nhập}

Ta thường đối mặt với lượng lớn dữ liệu, nhưng dữ liệu không hẳn là vô trật
tự. Những thành phố nối với nhau bằng đường tạo thành quan hệ đồ thị; quan hệ
cấp trên - cấp dưới trong cơ quan là quan hệ cây; xếp hạng điểm số ở trường
là quan hệ tuyến tính.

Đồ thị, cây và quan hệ tuyến tính là ba loại quan hệ thường gặp nhất trong
đời sống cũng như trong thi tin học. Dù thi tin học nhấn mạnh việc khai thác
và sử dụng đầy đủ thông tin đầu vào, đôi khi quan hệ quá phức tạp lại khiến
ta mất phương hướng. Tư tưởng chính ở đây là đơn giản hóa cấu trúc dữ liệu:
có thể biến đồ thị thành cây, biến cây thành cấu trúc tuyến tính. Trong quá
trình này chắc chắn mất đi một phần thông tin, nhưng các ví dụ dưới đây cho
thấy đó là một hướng tư duy khả thi.

\section{Đơn giản hóa quan hệ đồ thị}

\subsection{Bài toán đi thuyền}

Trường trung học Yali có \(n\) học sinh đi công viên chèo thuyền. Mỗi thuyền
tối đa chở hai người. Nếu hai học sinh cùng họ hoặc cùng tên, họ có thể ngồi
chung thuyền. Nhà trường muốn mọi học sinh đều lên thuyền, nhưng chi phí thuê
thuyền cao, nên cần thuê ít thuyền nhất.

Một cách tự nhiên là xem mỗi học sinh là một đỉnh. Nếu hai học sinh cùng họ,
nối họ bằng cạnh đỏ; nếu cùng tên, nối bằng cạnh xanh. Khi đó bài toán trở
thành tìm phủ cạnh nhỏ nhất, tương đương với tìm ghép cực đại. Nhưng đồ thị
không nhất thiết là hai phía, nên ghép cực đại cần thuật toán blossom, rất
phức tạp.

Giả sử trước hết đồ thị liên thông. Ta dùng một cây nhị phân để biểu diễn đồ
thị: con trái của mỗi nút, nếu tồn tại, có cùng họ với nút đó; con phải, nếu
tồn tại, có cùng tên.

Cách dựng cây:
\begin{enumerate}
  \item Giả sử cây hiện có \(k\) đỉnh. Nếu \(k=n\), kết thúc.
  \item Vì đồ thị liên thông, trong các đỉnh chưa vào cây có một đỉnh \(P\)
  nối với một đỉnh \(T\) đã ở trong cây.
  \item Nếu \(P\) và \(T\) cùng họ, đi liên tục theo con trái từ \(T\) đến khi
  gặp đỉnh \(X\) không có con trái. Vì \(X\) cùng họ với \(T\), \(X\) cũng
  cùng họ với \(P\); đặt \(P\) làm con trái của \(X\).
  \item Nếu \(P\) và \(T\) cùng tên, làm tương tự theo con phải.
\end{enumerate}

Như vậy luôn có thể tăng kích thước cây thêm \(1\) cho đến khi mọi đỉnh đều
nằm trong cây.

Sau khi có cây, xét một lá \(P\), cha của nó là \(F\). Nếu \(P\) là con duy
nhất của \(F\), cho \(P\) và \(F\) đi chung thuyền, rồi xóa cả hai khỏi cây.
Nếu \(F\) có hai con, gọi con còn lại là \(Q\). Không mất tính tổng quát, giả
sử \(P\) là con trái, \(Q\) là con phải.

Nếu \(F\) là gốc, còn ba học sinh chưa lên thuyền; ít nhất cần hai thuyền.
Cho \(F\) và \(P\) đi chung, còn \(Q\) đi một thuyền.

Nếu \(F\) không phải gốc, xét cha \(F'\) của nó. Trường hợp \(F\) là con trái
của \(F'\): cho \(F\) và \(Q\) đi chung, xóa \(F,Q\) khỏi cây. Vì \(F'\) cùng
họ với \(F\), và \(F\) cùng họ với \(P\), nên \(F'\) cũng cùng họ với \(P\);
đặt \(P\) làm con trái của \(F'\), cây vẫn giữ được tính chất. Trường hợp
\(F\) là con phải làm tương tự với quan hệ cùng tên.

Lặp lại quá trình trên sẽ liên tục giảm kích thước cây. Với cây có \(n\) đỉnh,
số thuyền cuối cùng là
\[
  \left\lceil\frac n2\right\rceil,
\]
rõ ràng là tối ưu. Nếu đồ thị không liên thông, xử lý từng thành phần liên
thông rồi cộng kết quả.

Độ phức tạp của thuật toán là \(O(n)\), tốt hơn nhiều so với thuật toán ghép,
và độ phức tạp cài đặt cũng đơn giản hơn hẳn. Điều đáng chú ý là ta đã bỏ qua
nhiều cạnh trong đồ thị, nhưng quan hệ được giữ lại đủ để giải bài toán tối
ưu.

\subsection{Nhận diện đồ thị xương rồng}

Một đồ thị có hướng được gọi là \term{xương rồng} nếu:
\begin{itemize}
  \item nó liên thông mạnh;
  \item mỗi cạnh thuộc đúng một chu trình.
\end{itemize}

Cho đồ thị có \(n\) đỉnh và \(m\) cạnh, với \(1\le n,m\le 10^5\), cần kiểm tra
nó có phải đồ thị xương rồng hay không.

Một thuật toán dễ nghĩ là liên tục tìm một chu trình, kiểm tra giữa các đỉnh
không kề trong chu trình có cạnh phụ hay không; nếu không thì co chu trình
thành một đỉnh rồi tiếp tục. Nhưng chỉ riêng việc co đỉnh và cập nhật quan hệ
giữa chu trình, đỉnh và cạnh đã rất phức tạp.

Cách khác là chạy DFS và xét cây DFS. Nếu đồ thị là xương rồng, cây DFS của
nó có các tính chất sau.

\paragraph{Tính chất 1.}
Cây DFS của đồ thị xương rồng không có cạnh ngang. Thật vậy, nếu có cạnh
ngang từ một nhánh sang nhánh khác, do đồ thị liên thông mạnh, sẽ tồn tại
đường quay lại. Dù đường này giao hay không giao với đường tổ tiên tương ứng,
một số cạnh sẽ thuộc hai chu trình, mâu thuẫn với định nghĩa xương rồng.

\paragraph{Tính chất 2.}
Đặt \(\operatorname{DFS}(v)\) là thứ tự thăm của \(v\) trong DFS. Đặt
\(\operatorname{Low}(v)\) là giá trị DFS nhỏ nhất có thể đi tới bằng một cạnh
đi ra trực tiếp từ \(v\) hoặc từ các hậu duệ của \(v\). Với mọi cạnh cây
\((v,u)\), trong đó \(u\) là con của \(v\), phải có
\[
  \operatorname{Low}(u)\le \operatorname{DFS}(v).
\]
Nếu không, cạnh \((v,u)\) là cầu, điều không thể xảy ra trong đồ thị liên
thông mạnh.

\paragraph{Tính chất 3.}
Với mỗi đỉnh \(v\), gọi \(a(v)\) là số con \(u\) của \(v\) thỏa
\(\operatorname{Low}(u)<\operatorname{DFS}(v)\), và \(b(v)\) là số cạnh ngược
từ chính \(v\). Khi đó phải có
\[
  a(v)+b(v)<2.
\]
Nếu tổng này từ \(2\) trở lên, sẽ có cạnh thuộc hai chu trình.

Ba tính chất trên cho ta thuật toán DFS tuyến tính. Nếu cây DFS vi phạm bất
kỳ tính chất nào, đồ thị không phải xương rồng. Ngược lại, nếu thỏa tất cả,
có thể chứng minh bằng phân tích cạnh cây và cạnh ngược rằng đồ thị là xương
rồng. Độ phức tạp thời gian là \(O(n+m)\).

Điểm bản chất là ta mô tả lại đồ thị dưới dạng cây DFS. Với cây, quan hệ tổ
tiên - hậu duệ rõ ràng hơn, và các cạnh ngược, cạnh ngang có thể được xử lý
tách biệt. Hình thức đơn giản hơn làm quan hệ dữ liệu nổi lên rõ ràng.

\section{Đơn giản hóa quan hệ cây}

\subsection{Thống kê trên cây}

Cho một cây có \(n\) đỉnh, đánh số \(1,2,\ldots,n\). Với đỉnh \(v\), định
nghĩa \(t(v)\) là số hậu duệ của \(v\) có số hiệu nhỏ hơn \(v\). Cần in
\[
  t(1),t(2),\ldots,t(n).
\]

Thuật toán trực tiếp \(O(n^2)\) quá chậm. Ta có thể giải bằng một cách khác:
biến cây thành hai dãy.

Thực hiện DFS và ghi lại các đỉnh theo thứ tự thăm, thu được \term{dãy DFS}.
Sau đó đảo thứ tự các con của mỗi đỉnh rồi DFS lại, thu được \term{dãy DFS
ngược}. Với một đỉnh \(v\), phần sau \(v\) trong dãy DFS bao phủ một vùng,
phần sau \(v\) trong dãy DFS ngược bao phủ một vùng khác; giao của hai vùng
này chính là tập hậu duệ của \(v\), còn phần bị tính lệch liên quan đến chính
\(v\) và các tổ tiên trực tiếp của nó.

Định nghĩa \(f(v,S)\) là số đỉnh trong tập hoặc vùng \(S\) có số hiệu nhỏ hơn
\(v\). Khi đó theo nguyên lý bao hàm - loại trừ:
\[
\begin{aligned}
t(v)=&\ f(v,\text{phần sau \(v\) trong dãy DFS})\\
&+f(v,\text{phần sau \(v\) trong dãy DFS ngược})\\
&+f(v,\text{các tổ tiên trực tiếp của \(v\)})-(v-1).
\end{aligned}
\]

Số tổ tiên trực tiếp nhỏ hơn \(v\) có thể tính bằng một lần DFS kết hợp cây
đoạn. Với hai dãy DFS, số phần tử phía sau nhỏ hơn phần tử hiện tại có thể
tính bằng cây đoạn hoặc Fenwick. Vì vậy toàn bộ bài toán được giải trong
\[
  O(n\log n).
\]

Điểm tinh tế là cây được ánh xạ thành hai dãy. Quan hệ hậu duệ trong cây vốn
có tính phân tầng; sau biến đổi, nó trở thành quan hệ trước - sau trong dãy,
và nhờ đó có thể dùng cấu trúc dữ liệu tuyến tính quen thuộc.

\subsection{Tổ tiên chung gần nhất}

Cho một cây \(n\) đỉnh và nhiều truy vấn dạng: tổ tiên chung gần nhất của
\(p\) và \(q\) là gì?

Ta biểu diễn cây bằng một dãy Euler. Nếu cây chỉ có một đỉnh, dãy là
\((\mathit{Root})\). Nếu gốc có các cây con tương ứng với các dãy
\[
  L_1,L_2,\ldots,L_t,
\]
thì dãy của cây là
\[
  (\mathit{Root},L_1,\mathit{Root},L_2,\mathit{Root},\ldots,
  L_t,\mathit{Root}).
\]

Ghi kèm độ sâu của mỗi lần xuất hiện. Khi cần tìm LCA của \(p,q\), chọn một
lần xuất hiện bất kỳ của \(p\) và một lần xuất hiện bất kỳ của \(q\) trong
dãy. Trong đoạn giữa hai vị trí đó, đỉnh có độ sâu nhỏ nhất chính là tổ tiên
chung gần nhất.

Chứng minh cũng theo đệ quy. Giả sử kết luận đúng với mọi cây con
\(L_1,\ldots,L_t\). Nếu \(p,q\) nằm trong cùng một cây con \(L_i\), kết luận
theo giả thiết quy nạp. Nếu \(p\in L_i\), \(q\in L_j\) với \(i\ne j\), LCA
của chúng là \(\mathit{Root}\). Trong cách dựng dãy, giữa hai dãy con bất kỳ
đều có một lần xuất hiện của \(\mathit{Root}\), và gốc có độ sâu nhỏ nhất, nên
kết luận đúng.

Do đó bài toán LCA được đưa về bài toán RMQ trên dãy độ sâu. Có thể dùng cây
đoạn để trả lời mỗi truy vấn trong \(O(\log n)\), hoặc dùng RMQ tĩnh để trả
lời trong \(O(1)\).

\section{Kết luận}

Chủ đề của bài viết là đơn giản hóa quan hệ dữ liệu. Hai dạng đơn giản hóa
quan trọng là từ đồ thị sang cây, và từ cây sang dãy. Những ví dụ trên cho
thấy đây không phải là việc bỏ thông tin một cách tùy tiện, mà là thay đổi
hình thức biểu diễn để làm rõ mâu thuẫn cốt lõi của bài toán.

Nói nghiêm ngặt, ``đơn giản hóa'' chỉ là một biến đổi hình thức: dùng cấu trúc
đơn giản và nguyên thủy hơn để biểu diễn quan hệ vốn được thể hiện bằng cấu
trúc phức tạp. Sự đơn giản này thường làm quan hệ dữ liệu rõ hơn, tạo điều
kiện để nhìn thấy bản chất bài toán. Đặc biệt, dùng cây DFS và dãy DFS để xử
lý các bài toán đồ thị, cây là một tư tưởng rất thực dụng và hiệu quả.

Ngoài DFS còn có nhiều ứng dụng khác của việc đơn giản hóa quan hệ dữ liệu.
Chẳng hạn trong lý thuyết đồ thị, kiểm tra đồ thị dây cung bằng \term{thứ tự
loại bỏ hoàn hảo} cũng là một ví dụ tốt.

\section*{Tài liệu tham khảo}

\begin{itemize}
  \item Đề thi tỉnh Hồ Nam.
  \item Bài viết của Xiao Zhou, đội tuyển tập huấn quốc gia tin học năm 2000.
  \item \textit{Introduction to Algorithms}.
  \item Bài toán gốc của Lei Tao, Trường trung học Yali, Trường Sa, Hồ Nam.
\end{itemize}

\end{document}
